\documentclass[a4paper]{article}
\usepackage[pdftex]{graphicx}
\usepackage[utf8]{inputenc}
\usepackage{enumerate}
\usepackage{icomma}
\usepackage{amssymb}
\usepackage{tikz}
\usepackage{href-ul}
\hypersetup{
	colorlinks=true,
	linkcolor=blue,
	urlcolor=blue}
\usepackage{geometry}
\geometry{a4paper, top=15mm, left=15mm, right=15mm, bottom=15mm,
	headsep=10mm, footskip=12mm}

\begin{document}
\begin{tikzpicture}(10,3)
	\draw[ultra thick](2,0) --(10,0) -- (10,1.3) --(2,1.3) -- (2,0);
	\draw[fill=black](2,0)-- (0,.6) -- (2,.6) -- (2,0);
	\node at (6,.6) {\large Lösungsblatt von \href{https://www.okuyakl.de}{www.okuyakl.de}};
\end{tikzpicture}
\vspace{0.5 cm}

\noindent{\bf Aufgabe 1. a)}\\
Die relative Häufigkeit ist: 
$$h(3)={8\over 40} = { 1 \over 5}$$

\noindent{\bf Aufgabe 1. b)}\\
$$H(4)= 40 \cdot 0,125 = 5 $$
Die Augenzahl 4 wurde 5 mal geworfen.

\noindent{\bf Aufgabe 1. c)}\\
Die absolute Häufigkeit der 6 ist:
$$H(6)=40-(7+4+8+5+6)=10$$

\noindent{\bf Aufgabe 1. d)}\\
Wir addieren die Häufigkeiten:
$$H(X\le 3)=H(1)+H(2)+H(3)= 7+4+8=19$$

\noindent{\bf Aufgabe 2.}\\ 
\begin{itemize}
	\item Für $\Omega_1$ und $\Omega_3$ liegt ein Laplace-Experiment vor, weil jedes Teilereignis gleich wahrscheinlich ist.
	\item Für $\Omega_2$ liegt kein Laplace-Experiment vor, weil die Wahrscheinlichkeit für eine 6 nicht gleich der Wahrscheinlichkeit für keine 6 ist.
\end{itemize}

\noindent{\bf Aufgabe 3.}\\
\begin{itemize}
	\item $$p(A)={1 \over 6}$$
	\item $$B=\{2;4;6\} \qquad p(B)={1 \over 2}$$
	\item $$C=\{2;1\}  \qquad  p(C)={1 \over 3}$$
\end{itemize}

\noindent{\bf Aufgabe 4.}\\
Der Ergebnisraum ist:
$$\Omega=\{WW; MM; MW; WM\}$$ 
Alle Elemente besitzen die gleiche Wahrscheinlichkeit $p={1 \over 4}$. Die Ereignismenge E= ,,mindestens ein Junge" besteht aus drei Elementen, also ist die Wahrscheinlichkeit hierfür:$p(E)={3 \over 4}$

\noindent{\bf Aufgabe 5.}\\
Die Ereignismenge ist:
$$E=\{55; 56; 65; 66\}$$
Jedes Element hat die Wahrscheinlichkeit $p={1 \over 36}$, damit ist
$$p(E)={4 \over 36}={1 \over 9}$$


\noindent
\begin{minipage}{0.4\textwidth}
	\noindent{\bf Aufgabe 6. a)}\\
	\includegraphics[width=7 cm]{bauzw059}
\end{minipage}
\hfill
\begin{minipage}{0.6\textwidth}
	\noindent{\bf Aufgabe 6. b)}\\
	$$\Omega=\{WWW; WWZ; WZW; WZZ; ZWW; ZWZ; ZZW; ZZZ \}$$

\end{minipage}
	
\noindent{\bf Aufgabe 6. c}\\
\begin{itemize}
	\item $$A=\{WWW; ZZZ\}; \qquad p(A)= {1 \over 4}$$
	\item $$B=\{ZWW; WZW; WWZ\}; \qquad p(B)= {3 \over 8}$$
	\item $$C=\Omega- \{ZZZ\} \qquad p(C)= {7 \over 8}$$
	\item $$D=\{WWW; ZWZ; WWZ; ZWW\}; \qquad p(D)= {1 \over 2}$$
	\item $$A=\{ZWZ\}; \qquad p(E)= {1 \over 8}$$
\end{itemize}

\begin{center}
	\includegraphics[width=7 cm]{../../viecher/endcomic.pdf}
	
	Hier geht es zurück zum \href{https://www.okuyakl.de/math/m9zufL059/aa059.pdf}{Aufgabenblatt}
\end{center}

\end{document}

